\begin{definition}
    $\nlsp=\nsp{\log n}$
\end{definition}

Для класса \np у нас было 2 определения: через недетерминированные машины и через сертификаты. Для класса \nlsp тоже удобно иметь сертификатное определение. Однако если перенести определение механически, то получится не \nlsp, а снова \np: например, истинность 3-КНФ на некотором наборе можно проверить на логарифмической памяти. Нужно ограничить возможности машины по обращению к сертификату. Это делается так: машина может читать сертификат только один раз, слева направо.

\begin{theorem} 
    $A \in \nlsp \Longleftrightarrow \exists$ машина $V(x,s):$ сертификат $s$ полиномиальной длины от $|x|$ и читается 1 раз слева направо, также машина использует логарифмическую память на рабочей ленте, и $x \in A$.
\end{theorem}

\begin{proof}
    $\Rightarrow$ Пусть $A \in \nlsp$, т.е. распознаётся некоторой недетерминированной машиной $M$ на логарифмической памяти. В таком случае сертификат $s$ будем понимать как последовательность выборов для машины  $M$, а машина $V$ будет просто моделировать работу машины $M$ на соответствующей ветви и вернёт то же значение. Если $x \in A$, то есть и принимающая ветвь, именно её и возьмём в качестве сертификата. Если же $x \not\in A$, то все ветви отвергающие/есть не заканчивающиеся, значит, любой сертификат будет отвергнут.
    
    $\Leftarrow$ Если для языка есть логарифмический верификатор, то можно написать недетерминированную машину, которая будет угадывать биты сертификата один за другим по мере надобности. Она использует ровно такое же количество памяти и выдаст единицу, если сможет угадать верный сертификат.
\end{proof}

\begin{lemmanote}
    Заметим, что из-за того, что машина использует логарифмическую память, сертификат в этом определении автоматически получается полиномиальной длины: все последующие биты верификатор всё равно не успеет прочесть.
\end{lemmanote}
